\documentclass[11pt]{article}
\usepackage[margin=1in]{geometry}
\usepackage{amsmath,amssymb,amsthm,mathtools}
\usepackage[T1]{fontenc}\usepackage{lmodern,microtype,xurl}
\usepackage[colorlinks=true,linkcolor=blue,urlcolor=blue]{hyperref}
\newtheorem{theorem}{Theorem}[section]
\newtheorem{proposition}[theorem]{Proposition}
\newtheorem{lemma}[theorem]{Lemma}
\newtheorem{corollary}[theorem]{Corollary}
\theoremstyle{remark}\newtheorem{remark}[theorem]{Remark}
\newcommand{\R}{\mathbb R}\newcommand{\E}{\mathbb E}
\newcommand{\KL}{\operatorname{D}}\newcommand{\poly}{\operatorname{poly}}
\newcommand{\norm}[1]{\left\lVert#1\right\rVert}
\title{Structural Boolean circuits and subexponential sample bounds\\for the original Gaussian tanh posterior}
\author{ChatGPT-6 Astra \and Kaarel H\"anni \and Dmitry Vaintrob}
\date{26 October 2026}
\begin{document}\maketitle
\begin{abstract}
We give structural sufficient conditions for subexponential dependence
on learner depth in the original, unrestricted Gaussian tanh posterior.
A synchronous Boolean circuit with an early narrow layer can be evaluated
by retaining its prefix and replacing its suffix by a truth table. A
bounded-fanin circuit of small graph treewidth admits a shallow threshold
implementation by tree-decomposition dynamic programming. For planar
circuit graphs, a separator schedule improves the exponent to the square
root of graph size, apart from logarithmic terms. We account for
the size and coefficients of these constructions, and give the resulting
sample bounds at every realizable learner width, uniformly over canonical
Gaussian scales $0<t\le1$. The statistical step imports the established
uniform intrinsic-learning and scalar-continuation theorems; a
heterogeneous-layer prefix lemma needed here is proved explicitly.
An explicit unary-map compiler also gives polynomial sample bounds
for Boolean formulas of any fixed Horton--Strahler rank, including
short learner depths; rank at most one works at every depth $L\ge2$.
These results do not cover arbitrary polynomial-size deep Boolean circuits
or replace real tanh states by their signs.
\end{abstract}

\section{Model and statements}
Let $X_d=\{-1,1\}^d$, $f:X_d\to\{-1,1\}$, and $H=L-1\ge1$.
The learner has $H$ hidden tanh layers of common width $N$, input
$(1,x)$, and a scalar linear readout. All weights are independent
centered Gaussians, with entry variances
\[
 \operatorname{Var}(W^1_{ij})=t^2/d,\qquad
 \operatorname{Var}(W^j_{ik})=t^2/N\ (j\ge2),\qquad
 \operatorname{Var}(v_i)=t^2/N^2,
 \quad 0<t\le1.
\]
There are no additional hidden biases. A constant input coordinate may
be propagated as an ordinary neuron. Write $z_w(x)$ for the output and
\[
 A_f=\{w:f(x)z_w(x)>1/2\text{ for every }x\in X_d\},\qquad
 a_{N,t}(f)=-N^{-1}\log\mu_{N,t}(A_f).
\]
A width is \emph{realizable} if $\mu_{N,t}(A_f)>0$. The posterior
conditions on $|z_w(x)|>1/2$ at every cube input and on agreement with
all noiseless training labels. Training and test inputs have any common
distribution $\mathcal D$. Its confidence error is
\[
 \mathcal D\{x:\Pr_{\rm post}(\operatorname{sign}z_w(x)=f(x))\le2/3\}.
\]
Every statistical assertion below is at each fixed realizable width;
it does not assert one simultaneous event for infinitely many widths.

For $E\ge1$, abbreviate the sufficient sample bound
\begin{equation}\label{eq:samples}
 n\ge \frac C\varepsilon\left[
       e^{CE}(1+\log(1/\varepsilon))^C+\log(1/\delta)\right],
 \qquad 0<\varepsilon,\delta<1.
\end{equation}
Here and below $C$ denotes a universal constant that may increase.

\begin{theorem}[An early narrow cut]\label{thm:cut}
Suppose $f$ is computed by a synchronous fan-in-two Boolean circuit
with $s\ge1$ gates and depth $h\ge1$: inputs occur at level zero,
and every wire joins adjacent levels. Let $w_j$ be the number of
bits at level $j$, $1\le j\le h$. Put
\[
 K=\log(d+s+L+3),\qquad g=\frac K{\log(K+2)},\qquad
 E_{\rm cut}=K+\min_{1\le j\le h}(jg+w_j).
\]
If $L\ge C E_{\rm cut}$, then \eqref{eq:samples} with
$E=E_{\rm cut}$ bounds the confidence error by $\varepsilon$ with
training probability at least $1-\delta$, at every realizable width
and every $0<t\le1$. In particular one may replace $E_{\rm cut}$ by
\begin{equation}\label{eq:sqrt}
 E_{\rm size}=K+\sqrt{sK/\log(K+2)}.
\end{equation}
If inputs are also injected after a cut, the narrow-cut assertion holds
with $w_j$ replaced by $w_j+r_j$, where $r_j$ counts the distinct
input bits used strictly after that cut. The size-only bound
\eqref{eq:sqrt} requires the original synchronous, level-zero-input
hypothesis.
\end{theorem}

\begin{theorem}[Small treewidth]\label{thm:tree}
Suppose a fan-in-two Boolean circuit computing $f$ has $v\ge1$
vertices, including its input vertices, and its underlying undirected
graph has treewidth $k$. Set
\[
 E_{\rm tw}=\log(d+L+3)
       +(k+\log(v+d+L+3))\log(v+3).
\]
If $L\ge C E_{\rm tw}$, then \eqref{eq:samples} with $E=E_{\rm tw}$
has the same confidence-error, scale, and all-width conclusions.
\end{theorem}

\begin{theorem}[Planar and excluded-minor circuit graphs]\label{thm:planar}
Suppose the underlying undirected graph of a fan-in-two Boolean circuit
for $f$ is planar, with $v\ge1$ vertices including its input vertices.
Set
\[
 E_{\rm sep}=\sqrt v+\log(d+L+3)
                 +\log(v+3)\log(v+d+L+3).
\]
If $L\ge C E_{\rm sep}$, then \eqref{eq:samples} with $E=E_{\rm sep}$
has the same confidence-error, scale, and all-width conclusions.
The same statement holds for any fixed excluded-minor graph class,
with constants depending only on the excluded minor.
\end{theorem}

\begin{theorem}[Formula register rank]\label{thm:strahler}
Suppose $f$ has a fan-in-at-most-two Boolean formula with $s\ge1$
nodes and Horton--Strahler rank $r$. A leaf has rank zero, a unary
gate retains its child's rank, and a binary gate of child ranks
$a,b$ has rank $\max\{a,b\}$ if $a\ne b$ and rank $a+1$ otherwise.
Set
\[
 E_r=\log(d+L+3)+(r+1)\log(d+s+L+3).
\]
If $L\ge C E_r$, then \eqref{eq:samples} with $E=E_r$ has the same
confidence-error, scale, and all-width conclusions. In particular,
the sample bound is polynomial in $d+s+L$ at every fixed rank,
independently of the formula's original depth. Repeated occurrences
of the same input variable are allowed.
\end{theorem}

\begin{corollary}[Fixed-rank formulas at short learner depths]\label{cor:rankshort}
For the formula in Theorem~\ref{thm:strahler}, put $R=d+s+L+3$.
There are universal constants such that, for every $L\ge C(r+1)$,
a sufficient sample bound is
\[
 n\ge\frac C\varepsilon\left[
 R^{C(r+1)^2}(1+\log(1/\varepsilon))^C+\log(1/\delta)\right].
\]
This has the same confidence-error conclusion, at every realizable
width and every $0<t\le1$. If $r\le1$, a polynomial bound in
$d+s+L$ holds at every $L\ge2$.
\end{corollary}

For families with $d,s,v\le L^{O(1)}$, these theorems give an
$\exp(o(L))$ sample bound at fixed $\varepsilon,\delta$ whenever,
respectively,
\[
 \min_j(j\log L/\log\log L+w_j)=o(L),\quad
 s=o(L^2\log\log L/\log L),\quad\text{or}\quad
 k=o(L/\log L).
\]
The planar or fixed-excluded-minor bound is subexponential when
$v=o(L^2)$.
The displayed padding conditions then hold eventually. In particular,
they must not be omitted from finite-parameter versions of the results.
The conclusions restrict the \emph{target's} Boolean representation;
they impose no sparsity, treewidth, or architecture restriction on the
prior or posterior.

\section{The analytic learning inputs and a layerwise prefix}
The structural arguments below construct threshold circuits, charging
COPY gates and constant rails. The prefix lemma turns their layerwise
resources into Gaussian target witnesses. The following imported results
supply the posterior comparison; a VC argument covers smaller widths.
Let
\[
 b=t^{-2}-1,\qquad c_t=(t^{-2}-1+2\log t)/2,
\]
and let $V_*$ denote the critical scalar Bellman baseline of
\cite{critical,uniform}. Its numerical value is not needed.

\begin{proposition}[Imported intrinsic and scalar inputs]\label{prop:inputs}
The following statements hold for the model above.
\begin{enumerate}
\item For the actual target cost $a=a_{N,t}(f)$, or a deterministic
upper bound for it, define
\[
 E_a=(a+4/(Ht^2)-Hc_t-V_*+b)_+,
 \qquad J_a=\min\{\sqrt{HE_a},E_a/b\},
\]
omitting the second term when $b=0$. A sufficient sample size is
polynomial in $d,L,1+t^2a$, multiplied by $\exp(CJ_a)$ and the
statistical factors in \eqref{eq:samples}. The polynomial degrees
and constants are universal, and the conclusion holds at every
realizable width uniformly in $0<t\le1$.
\item Suppose a Boolean circuit prefix of $q$ hidden stages supplies
a target component $p\ge e^{-CS}$, diagonal at most one, critical
cost at most $1/(8L)$, and finite-width relative score and residual
accuracy $e^{-C'S}$ at widths $N\ge e^{C''S}$, for any required
fixed exponents after increasing constants. Suppose $S\ge
1+\log(d+L+2)+q$. Its Gaussian mean shifts at scale $t$ have cost
at most $1/(8Lt^2)$. If $H-q\ge CS$, the critical scalar program,
reused at scale $t$, gives the actual target-mass bound
\begin{equation}\label{eq:target}
 a_{N,t}(f)\le Hc_t+V_*+C(bS+S^2/H)
 \quad(N\ge e^{CS}),
\end{equation}
with constants and width threshold independent of $t$.
\end{enumerate}
\end{proposition}

The first item is the baseline-subtracted intrinsic bound in
\cite[Proposition 2.1, equation (4)]{uniform}, proved in the companion
note \texttt{uniform-sharp-intrinsic.md}. The second is the scalar
continuation and finite-width lift in
\cite[Proposition 3.3]{uniform} (label
\texttt{prop:long}), with the prefix supplied separately. Its proof
uses the prefix only through the listed seed, cost, and accuracy
bounds. The underlying finite-width circuit prefix is
\cite[Lemma \texttt{lem:critical-finite-circuit-prefix}]{critical}.
These are imported analytic theorems, not assertions reproved by the
combinatorial arguments below.

\begin{lemma}[A heterogeneous-layer prefix]\label{lem:prefix}
Consider a layered threshold circuit with $q\ge1$ stages, widths
$C_j\ge1$, a constant rail, unit normalized gate margins, and
row Euclidean norms at most $B_j\ge1$ at stage $j$. Set
\begin{equation}\label{eq:S}
 S=1+\log(d+L+3)+q\log(q+L+3)
       +\sum_{j=1}^q\log((C_j+1)(B_j+1)).
\end{equation}
For $L\ge q+1$ this circuit supplies the prefix hypotheses in
Proposition~\ref{prop:inputs}, with this $S$. Consequently, when
$L\ge CS$, \eqref{eq:samples} holds with $E=S$ uniformly over
all canonical scales and all realizable widths.
\end{lemma}
\begin{proof}
We give the resource dependence, including the finite-width
construction, to avoid replacing the sum in \eqref{eq:S} by a
maximum width times depth. Put $p_0=1/d$, and choose a common
gain $a$ and tolerance $\eta=8e^{4-a}$. Define
\[
 p_j=\frac{p_{j-1}}{8C_jqLa^2B_j^2},\qquad
 k_j=\lceil Np_j\rceil.
\]
For each wire prescribe a disjoint group of $k_j$ neurons. In the
first layer use mean $aw_{1,i}$ on a controlled row. At stage $j>1$,
put mean $aw_{j,i,g}/k_{j-1}$ on each coordinate in the preceding
group for wire $g$, and zero elsewhere. All row covariances retain
their original Gaussian values. The preactivation mean is exactly
$a\sum_g w_{j,i,g}m_{j-1,g}(x)$, where $m_{j-1,g}$ is that group's
empirical mean.

If $Np_q\ge1$, then $k_j\le2Np_j$, the groups fit, and the entropy
per controlled row is at most $a^2B_j^2/(2t^2p_{j-1})$. Thus
\[
 \frac1N\KL(P\Vert\mu^{\rm prefix}_{N,t})
 \le\sum_{j=1}^q
       \frac{2C_jp_j a^2B_j^2}{2t^2p_{j-1}}
 \le\frac1{8Lt^2}.
\]
Choose $\eta\le(4\max_j B_j\sqrt{W_j})^{-1}$, where
$W_1=d+1$ and $W_j=C_{j-1}$ thereafter. An incoming empirical
error at most $\eta$ then leaves preactivation mean of the correct
sign and magnitude at least $3a/4$. Gaussian variance is at most
two at the first stage and one thereafter. Gaussian exponential
moments and $1-\tanh u\le2e^{-2u}$ show that its tanh expectation
is within $\eta/4$ of the correct bit, and its second moment is
within $\eta/2$ of one.

Conditional Hoeffding bounds, united over all cube inputs, groups,
and stages, now give simultaneous gate-mean error at most $\eta$
provided
\[
 Np_q\ge C\eta^{-2}
       \left[d+\log\frac{q(1+\max_j C_j)}{\zeta}\right].
\]
The diagonal follows the scalar recurrence
$\bar q_j=(1-\theta_j)F_t(\bar q_{j-1})+\theta_j$,
where $\theta_j=C_j k_j/N$, $\bar q_0=(d+1)/d$, and
$F_t(u)=\E\tanh^2(t\sqrt uG)$.
Since $F_t$ is one-Lipschitz for $t\le1$, bounded-square
concentration gives diagonal error at most $2j\eta$.
The final target-group indicator, normalized to a unit vector,
therefore has score within $\sqrt p\eta$ of $\sqrt p f$, where
$p=k_q/N$, and residual-variance error at most $(2q+3)\eta$.
After the first stage the ideal diagonal is at most one. This is
the same score/residual seed used by the scalar continuation.

Take $a=A S$ for a sufficiently large universal $A$. Because
$S\ge q\log(q+L+3)$ and $S\ge q$, one has
\[
 q\log S=q\log q+q\log(S/q)\le2S.
\]
It follows that
\[
 \log(1/p_q)=\log d+
   \sum_j\log(8C_jqLa^2B_j^2)\le C(1+\log A)S.
\]
The gate-margin tolerance and any fixed seed-relative accuracy
requirement hold by increasing $A$: $\log(1/\eta)$ grows linearly
in $A$, while the last bound grows only logarithmically in $A$.
The concentration requirement is then $N\ge e^{CS}$ for
$\zeta\ge e^{-C'S}$, with arbitrary fixed $C'$. The construction
and all concentration estimates are uniform for $t\le1$.

For $L\ge CS$, Proposition~\ref{prop:inputs} gives
\eqref{eq:target}. In its intrinsic bound, $E_a\le
C(bS+S^2/H)$, absorbing $b$ and $4/(Ht^2)$ because
$S\ge1+\log(H+2)$. If $b\ge S/H$ use $E_a/b$; otherwise use
$\sqrt{HE_a}$. In either case $J_a=O(S)$. Also
$t^2a\le C(H+S+S^2/H)$, so all polynomial factors are absorbed
by $e^{CS}$.

This proves the conclusion for $N\ge e^{CS}$. For smaller widths,
the entire depth-$L$, width-$N$ tanh sign class has VC dimension
polynomial in its parameter and computation counts
\cite{km}. These counts are $e^{O(S)}$, since $d,L,N\le e^{O(S)}$.
The standard realizable VC consistency bound therefore makes every
posterior-supported function's test error at most $\varepsilon/3$
under \eqref{eq:samples}. Averaging over that posterior and applying
Markov's inequality gives confidence error at most $\varepsilon$:
at a point counted as an error, posterior mistake probability is at
least $1/3$. This handles every smaller realizable width without
requiring the compiled circuit to fit inside that width.
\end{proof}

\section{A prefix followed by one lookup}
We first record the resource accounting behind cone blocking.
Take $b$ consecutive fan-in-two Boolean levels. Each output depends
on at most $2^b$ entrance bits. Its truth table has a DNF with at
most $2^{2^b}$ terms, hence two threshold stages, whose widths and
normalized row norms have logarithms $O(2^b+\log(s+2))$.
Copy gates and a constant rail obey the same bounds.
Thus a depth-$j$ prefix has a threshold implementation with
\[
 q_0=O(1+j/b),\qquad
 \sum_{\rm prefix}\log((C_i+1)(B_i+1))
       =O((1+j/b)(K+2^b)).
\]
The same bounds hold if finitely many specified input bits must be
retained for later use. Choose $b$ within a constant factor of
$\log(K+2)$, truncated to $1\le b\le j$. Then $2^b=O(K)$ and
\begin{equation}\label{eq:prefixbudget}
 S_{\rm prefix}=O(K+jK/\log(K+2)).
\end{equation}
This is also the resource-sensitive blocking construction in
\cite[\texttt{thm:critical-boolean-blocking}]{critical}.

At a cut with $w$ bits, the whole suffix is a Boolean function of
those bits. Its DNF implementation has two threshold stages,
width $2^{O(w)}$, and normalized row norms $2^{O(w)}$.
It contributes $O(w+1)$ to the sum in \eqref{eq:S}, not
$O(q_0w)$. Combining with \eqref{eq:prefixbudget} gives
$S=O(K+jg+w)$ and proves the first part of
Theorem~\ref{thm:cut} by Lemma~\ref{lem:prefix}.
If some input bits enter after the cut, retain precisely those
distinct bits alongside the cut state. This proves the asserted
$w+r$ correction.

For completeness, first prune all gates that are not ancestors of
the designated output; this can only reduce the widths and gate count.
Take the first remaining level with at most $y$ gates.
It exists because the final output is one bit. All previous
levels have more than $y$ gates, so its index is at most
$s/y+1$. Therefore
\[
 E_{\rm cut}\le C\{K+y+(s/y)g\}.
\]
Choosing a positive integer $y$ within one of $\sqrt{sg}$ gives
\eqref{eq:sqrt}. This keeps the prefix of the synchronous
depth-reduction argument of G\'al--Jang \cite{gj}, instead of
using only its minimum-depth conclusion and losing size control.

\begin{remark}[Ordinary neural realizations]
The same threshold implementations have unit-margin tanh
realizations with logarithmic Frobenius product bounded by their
parameter $S$. Indeed, use a common output-bit tolerance
$\eta=e^{-CS}$ and tanh gain $a=O(S)$ at each threshold stage.
Gate scores retain a fixed sign gap, and saturation achieves that
tolerance. The layer norm is at most $aB_j\sqrt{C_j}$; hence
\[
 \log\prod_j\max\{1,\norm{W_j}_F\}
 \le q\log a+\sum_j\log B_j+\tfrac12\sum_j\log C_j
 =O(S).
\]
A final scalar multiplier gives unit output margin. This product
statement is additional evidence for the resource accounting;
the posterior conclusion above uses the Gaussian seed construction
and does not assume a generic product-to-target-mass theorem.
\end{remark}

\section{Tree-decomposition dynamic programming}
Start with a tree decomposition of the underlying circuit graph
of width $k$. The balanced-decomposition theorem gives a binary
decomposition of width at most $3k+2$, depth $O(\log(v+3))$, and
polynomially many bags \cite[Lemma 1]{js}. This external graph
theorem is the only structural input in this section.

Enlarge each bag by adding the at most two predecessors of every
circuit vertex already in it. Its size grows by a factor at most
three. This remains a tree decomposition: for any vertex $u$,
its new bags are its old connected bag subtree together with the
old subtrees of its successors. Each such subtree intersects the
old subtree of $u$, because a bag covers that circuit edge.
Their union is connected. Moreover each gate's full constraint
scope---the gate and its at most two inputs---is now contained
in a bag.

Introduce one Boolean variable $z_u$ for every circuit vertex.
For a fixed external input $x$, require every gate equation,
every input equation $z_u=x_i$, and $z_{\rm out}=1$.
Anchor each constraint at one bag containing its variables.
Because the circuit is acyclic, the whole system is satisfiable
exactly when the circuit accepts $x$.

For bag $B_t$ and assignment $a\in\{0,1\}^{B_t}$, let
$A_t(a;x)$ assert that $a$ extends to a satisfying assignment
of all constraints anchored below $t$. If $t$ has children
$\ell,r$, then
\begin{align}\label{eq:dp}
 A_t(a;x)={}&\operatorname{Local}_t(a;x)\ \wedge\notag\\[-2mm]
 &\bigvee_{b:\,b|_{B_t\cap B_\ell}=a|_{B_t\cap B_\ell}}
                  A_\ell(b;x)\ \wedge\notag\\[-2mm]
 &\bigvee_{c:\,c|_{B_t\cap B_r}=a|_{B_t\cap B_r}}
                  A_r(c;x).
\end{align}
For a missing child omit its factor. Running intersection shows
that the two child subtrees overlap only in variables already
fixed at the parent; hence compatible extensions glue. This
proves both directions of \eqref{eq:dp}.
For fixed $a$, the anchored gate constraints and output
condition are constants, while the input constraints are a
conjunction of at most $O(k+1)$ literals of $x$.

Compute all these local terms first. Retain them by copy gates
until their merge level, and evaluate \eqref{eq:dp} upward.
Each merge uses an OR stage and an AND stage. Finally OR over
all root assignments. There are $2^{O(k)}\poly(v+d)$ table
entries and connections, and $O(\log(v+3))$ threshold stages.
All thresholds can be chosen integral and normalized to margin
one; their row norms are at most $2^{O(k)}\poly(v+d)$.
Constant rails, copying pending tables, and retaining local
terms obey the same resource bounds. In particular the potentially
exponential table size is charged explicitly.

Substituting these resources into \eqref{eq:S} gives
\[
 S\le C\{\log(d+L+3)
          +(k+\log(v+d+L+3))\log(v+3)\}
       =O(E_{\rm tw}).
\]
Lemma~\ref{lem:prefix} proves Theorem~\ref{thm:tree}.
The preceding remark also gives a constant-margin tanh realization
of physical depth $O(\log(v+3))$ and logarithmic Frobenius product
$O((k+\log(v+d+3))\log(v+3))$, choosing its gain without the
statistical learner-depth parameter.

\section{Geometrically shrinking separators}
The weighted planar separator theorem \cite{lt} partitions any
$n$-vertex planar graph with nonnegative vertex weights into two
nonadjacent sides and a separator of at most $c\sqrt n$ vertices,
each side carrying at most $\beta=2/3$ of the total weight. We use
this theorem twice at each recursion step, for two different weights.

A subproblem consists of disjoint interior and boundary sets $I,B$,
with the induced graph on $I\cup B$; no interior vertex has an edge
outside this subproblem. Initially $I$ is the whole graph and
$B=\varnothing$. Write $m=|I|$, $b=|B|$. If $I=\varnothing$,
use the leaf bag $B$ and stop. Otherwise, first separate with weight
one on $I$ and zero on $B$, obtaining separator $Z_0$ and two
sides $U_0,U_1$, each with at most $\beta m$ interior vertices.
Within each nonempty $U_i$, apply the weighted theorem with weight
one on $B\cap U_i$ and zero elsewhere. If this weight is zero,
retain $U_i$ as one side without a separator. Let $Z_1,Z_2$ be
the additional separators and $P_a$ the at most four resulting
sides. Put
\[
 Z=Z_0\cup Z_1\cup Z_2,\qquad
 \mathcal B_t=B\cup Z,\qquad
 I_a=I\cap P_a,\quad B_a=(B\cap P_a)\cup Z.
\]
The children are the induced subproblems on $I_a\cup B_a$.
Duplicating all of $Z$ into each child is allowed. Their sizes obey
\begin{equation}\label{eq:boundary}
 |I_a|\le\beta m,\qquad
 |B_a|\le\beta b+3c\sqrt{m+b}.
\end{equation}
Different $P_a$ have no edges between them, so the boundary invariant
persists. The bags $\mathcal B_t$ form a tree decomposition:
edges are either covered in a parent bag or passed to one child,
and a vertex appearing in several children belongs to the parent
bag. Induction gives edge coverage and running intersection.

Put $M_\ell=\beta^\ell v$. Choose a universal $A$ large enough that
\[
 \beta A+3c\sqrt{1+A}\le A\sqrt\beta;
\]
such an $A$ exists because $\beta<\sqrt\beta$.
At depth $\ell$, interior size is at most $M_\ell$ and boundary
size is at most $A\sqrt{M_\ell}$. Indeed, when $M_\ell\ge1$,
\eqref{eq:boundary} and this induction hypothesis give
\begin{align*}
 |B_a|&\le\beta A\sqrt{M_\ell}
                 +3c\sqrt{M_\ell+A\sqrt{M_\ell}}\\
 &\le(\beta A+3c\sqrt{1+A})\sqrt{M_\ell}
 \le A\sqrt{M_{\ell+1}}.
\end{align*}
When $M_\ell<1$, interiors are empty and recursion has stopped.
Thus the decomposition has depth $O(\log(v+3))$, polynomially many
bags, and every bag at depth $\ell$ has size
$O(\sqrt{M_\ell})$. The last level has constant bag size.
Enlarge bags by gate predecessors exactly as in the preceding
section. Their sizes increase by at most three, while all gate
scopes become covered; no planarity of the enlarged scope graph
is needed.

Now perform the same bag-assignment dynamic program, scheduling
levels from deepest to root. Each table entry is a local term AND
at most four compatible-child ORs, so each level uses only a
constant number of threshold stages. Carry the original Boolean
input bits and a constant rail throughout, and compute each bag's
local literal constraints only when that bag's level is reached.
Leaf tables are also created at their own level. Child tables are
discarded after the next merge. In particular, large root tables
are not created at the first physical layer and retained throughout;
doing so would lose the following improvement.

At level $\ell$, logarithmic widths and normalized row norms are
\[
 O(\sqrt{M_\ell}+\log(v+d+3)).
\]
This includes input rails, local terms, all tables at that level,
and their connections. The geometric series satisfies
\[
 \sum_{\ell\ge0}\sqrt{M_\ell}
       \le\frac{\sqrt v}{1-\sqrt\beta}.
\]
With $q=O(\log(v+3))$ physical threshold stages, the sum in
\eqref{eq:S}, including its $q\log(q+L+3)$ term, is therefore
$O(E_{\rm sep})$. Lemma~\ref{lem:prefix} proves the planar case of
Theorem~\ref{thm:planar}.

For a fixed excluded minor, use the weighted
$O(\sqrt n)$ separator theorem of Alon--Seymour--Thomas
\cite{ast}; the same argument applies with its fixed constant.
Its component balance can be grouped into two sides of weight at
most $2/3$. This proves the extension.

\section{Flattening formula paths by their unary transformation monoid}
We prove a threshold implementation of a size-$s$, rank-$r$ formula
with depth $O(r+1)$, width polynomial in $d+s+2$, and unit-margin
integer gate coefficients of polynomial magnitude. Its polynomial
degrees are independent of $r$.

First consider $m$ unary Boolean maps, writing bits as $0,1$.
The map $u_i$ is specified by $a_i=u_i(0)$ and $b_i=u_i(1)$.
Put $R_i=\boldsymbol1\{a_i=b_i\}$ and
$F_i=\boldsymbol1\{a_i=1,b_i=0\}$. When $R_i=1$, the map
resets its argument to $a_i$; otherwise it is either identity
or a flip, according to $F_i$. Thus the composite applied to
$y_0$ equals $y_0$ XOR the parity of all flips if there is no
reset. If $j$ is the last reset, it equals $a_j$ XOR the
parity of flips with indices greater than $j$.

All these suffix parities have constant-depth threshold circuits.
For $0\le j\le m$ and $1\le k\le m-j$, form
\[
 U_{j,k}=\boldsymbol1\left\{\sum_{i>j}F_i\ge k\right\}.
\]
The alternating sum
$\sum_{k=1}^{m-j}(-1)^{k+1}U_{j,k}$ is exactly the suffix
parity; a second-stage threshold at $1/2$ gives its bit $P_j$.
For the empty suffix put $P_m=0$.
The count scores $2\sum_{i>j}F_i-(2k-1)$ and the parity score
$2\sum_k(-1)^{k+1}U_{j,k}-1$ have integer gap at least one.
There are $O(m^2)$ gates and all coefficient magnitudes and row
norms are polynomial in $m+1$.

The mutually exclusive last-reset indicators are
\[
 J_j=R_j\wedge\bigwedge_{i>j}\neg R_i\quad(1\le j\le m),
 \qquad J_0=\bigwedge_{i\ge1}\neg R_i.
\]
Put $a_0=y_0$. The answer is
\[
 \bigvee_{j=0}^m\bigl[J_j\wedge(a_j\mathbin{\mathrm{XOR}}P_j)\bigr].
\]
Each XOR has a fixed-size constant-depth Boolean implementation.
The indicators and final OR are threshold gates. This explicitly
computes composition in the four-element monoid consisting of
identity, flip, reset to zero, and reset to one, in constant
threshold depth and $O((m+1)^2)$ gates. A constant rail supplies
biases. Conversion from $0,1$ to signed Boolean wires preserves
the stated polynomial resources and unit margins.

Now form maximal formula paths using parent--child edges whose
endpoints have equal rank. At a rank-$k$ binary node there is at
most one child of rank $k$, so these components really are paths.
Every off-path child has smaller rank. At the bottom of a path,
either a leaf supplies the starting bit, or a gate whose children
have smaller rank computes it in constant depth. Going upward,
each gate acts as a unary map on the carried bit, with an already
computed off-path bit as its other argument. For a fixed gate
truth table, $a_i,b_i$ are unary functions of that off-path bit,
so they, $R_i$, and $F_i$ take only constant depth to prepare.
The preceding compiler evaluates the entire path in constant depth.

Process ranks $0,1,\ldots,r$, all paths of one rank in parallel.
Keep input bits, the constant rail, and computed path-root values
on ordinary copy wires until needed. There are at most $s$ such
values, and the sum of squared path lengths is at most $s^2$.
Consequently each rank uses polynomial width, connections, and
coefficient row norms in $d+s+2$, with universal polynomial degrees.
The total threshold depth is $O(r+1)$. Substitution into
\eqref{eq:S} gives $S=O(E_r)$; Lemma~\ref{lem:prefix} proves
Theorem~\ref{thm:strahler}.

The single-bit path is essential: this is not a constant-depth
compiler for arbitrary constant-width branching programs, whose
state transformation monoids can be non-solvable. A balanced
binary formula can have $r=\Theta(\log s)$, so the general formula
bound remains quasipolynomial. The argument reduces target depth
semantically before the Gaussian prefix construction; it does not
discard the costs of its physical neuron groups.

\subsection{Short padding and the rank-one Fourier seed}
The imported uniform threshold theorem
\cite[Theorem 1.1]{uniform} includes a short-padding assertion:
a threshold target of $q$ stages, width $C_0$, and normalized
coefficient bound $B$ has sufficient sample resource
\[
 (d+C_0+B+q+L+4)^{C(q+1)^2}
\]
at every $L\ge q+1$, all realizable widths, and all $0<t\le1$.
Apply this to the formula compiler, where $q\le C(r+1)$ and
$C_0,B$ are bounded by a universal polynomial in $d+s+2$.
This proves the first assertion of Corollary~\ref{cor:rankshort}.
Unlike the sharper linear-in-$r$ exponent, this statement does
not need logarithmic learner-depth padding.

For the remaining assertion, define the Fourier variation
$A(g)=\sum_{T\subseteq[d]}|\widehat g(T)|$.
It obeys $A(gh)\le A(g)A(h)$ on the Boolean cube, even when
$g,h$ depend on overlapping variables. A rank-zero formula is
a literal or constant, with variation one. In a rank-one formula,
all rank-one nodes lie on one path; each off-path value is a
literal or constant. The signed starting bit $f_0=2y_0-1$
depends on at most two literals, so $A(f_0)\le2$ by Parseval
and Cauchy--Schwarz.

In the notation above, $R_i$ and $1-R_i$ are literal indicators
or constant bits and have variation at most one. The signed
reset value $g_i=2a_i-1$ and flip sign $\eta_i=(-1)^{F_i}$
are signed literals or constants and have variation one.
The exact signed output is
\[
 f=J_0 f_0\prod_{i=1}^m\eta_i
          +\sum_{j=1}^m J_jg_j\prod_{i>j}\eta_i,
\]
where $J_0=\prod_i(1-R_i)$ and
$J_j=R_j\prod_{i>j}(1-R_i)$.
Submultiplicativity and the triangle inequality give
$A(f)\le2+m\le s+2$.
The uniform Fourier-target theorem
\cite[Theorem 1.2]{uniform} supplies a sample bound polynomial
in $d+L+A(f)+2$ at every $L\ge2$, every realizable width, and
every $0<t\le1$, with the usual statistical factors.
This proves the rank-one assertion. It uses the Fourier seed and
the imported theorem's finite-width completion; the threshold
compiler need not fit in that learner's width or hidden depth.

\section{Scope of the conclusions}
These are positive target classes for the original posterior,
including its unrestricted background neurons and all canonical
Gaussian scales. Small compiled width is unnecessary at an
individual learner width: at the smaller widths, the proof uses
the full finite-width class rather than forcing the compilation
into the learner. Realizability remains necessary for conditioning
on noiseless target labels.

The bounds do not solve the generic Boolean-depth question.
A polynomial-size circuit can have large treewidth and no early
sublinear-bit cut, and need not lie in a fixed excluded-minor class.
Generic planarization can have quadratic size cost, which removes
the sublinear-in-size gain of Theorem~\ref{thm:planar}.
Size-only depth reductions need not preserve
the resource sum \eqref{eq:S}. Moreover a bounded-width tanh
teacher stores real-valued states; its hidden signs need not be
sufficient statistics for the next layer. Applying these results
to such a teacher requires an independently justified Boolean
simulation with the stated structural resources. No cutset or
treewidth limitation of one representation is an intrinsic neural
lower bound, and no exponential sample lower bound is proved here.

\begin{thebibliography}{9}
\bibitem{uniform} ChatGPT-6 Astra, Kaarel H\"anni, and Dmitry Vaintrob.
\emph{Uniform prior-scale bounds for structured Boolean targets}.
Companion to \emph{Depth dependence, parallel circuits, and Boolean
learning}, 2026. Repository package
\path{works-in-progress/boolean-depth-barriers/}; source
\path{source/uniform-target-classes.tex}, with the six supporting
\path{proof-notes/} files identified in its proof-dependency section.

\bibitem{critical} ChatGPT-6 Astra, Kaarel H\"anni, and Dmitry Vaintrob.
\emph{Critical-depth analysis for the Gaussian tanh posterior}, 2026.
Repository package \path{works-in-progress/critical-depth/}; in
particular \path{source/critical-finite-width-circuit-prefix.tex},
\path{source/critical-optimized-circuit-blocking.tex}, and the
scalar-steering, finite-channel, and intrinsic-learning modules.

\bibitem{gj} Anna G\'al and Jing-Tang Jang.
\emph{The Size and Depth of Layered Boolean Circuits}.
LATIN 2010; journal version, Information Processing Letters, 2010.
\url{https://www.cs.utexas.edu/~panni/sizedepth.pdf}.
\url{https://doi.org/10.1016/j.ipl.2010.11.023}.

\bibitem{js} Maurice Jansen and Jayalal Sarma.
\emph{Balancing Bounded Treewidth Circuits}. Lemma 1, citing
Bodlaender's balanced-tree-decomposition theorem.
\url{https://www.cse.iitm.ac.in/~jayalal/research/papers/JS10.pdf}.

\bibitem{lt} Richard J. Lipton and Robert Endre Tarjan.
A separator theorem for planar graphs.
\emph{SIAM Journal on Applied Mathematics} 36(2) (1979), 177--189.
\url{https://doi.org/10.1137/0136016}.
Original report:
\url{https://i.stanford.edu/pub/cstr/reports/cs/tr/77/627/CS-TR-77-627.pdf}.

\bibitem{ast} Noga Alon, Paul Seymour, and Robin Thomas.
A separator theorem for graphs with an excluded minor and its applications.
\emph{STOC}, 1990.
\url{https://research.ibm.com/publications/separator-theorem-for-graphs-with-an-excluded-minor-and-its-applications}.

\bibitem{km} Marek Karpinski and Angus Macintyre.
Polynomial bounds for VC dimension of sigmoidal and general
Pfaffian neural networks. \emph{Journal of Computer and System
Sciences} 54 (1997), 169--176.
\url{https://doi.org/10.1006/jcss.1997.1477}.
\end{thebibliography}
\end{document}
